(1) 设复波数 $k=k_r+ik_i$，取相位约定 $e^{i(kx-\omega t)}$。代回波动方程得到
\[
-Tk^2=-\mu\omega^2-i\beta\omega,\quad
k^2=\frac{\mu\omega^2+i\beta\omega}{T}.
\]
实际弦位移由实部构成，因此 $e^{ikx}=e^{ik_rx-k_ix}$。满足向右传播而不在 $x\to+\infty$ 发散的物理根必须有 $k_r>0,k_i>0$；另一平方根会产生空间增长而不满足半无限弦的辐射条件。

(2) 令 $A_0=\mu\omega^2/T>0,\ B_0=\beta\omega/T>0$，联立 $k_r^2-k_i^2=A_0$、$2k_rk_i=B_0$ 得
\[
k_r=\sqrt{\frac{\sqrt{A_0^2+B_0^2}+A_0}{2}},\qquad
k_i=\sqrt{\frac{\sqrt{A_0^2+B_0^2}-A_0}{2}}.
\]
故稳定波形为
\[
y(x,t)=A e^{-k_ix}\cos(k_rx-\omega t).
\]
$1/k_i$ 是振幅衰减的 $1/e$ 长度，空间相位波长为 $2\pi/k_r$；两者都随 $\omega,\beta$ 改变。

(3) 将 PDE 乘以 $y_t$ 并按乘积法则整理：
\[
\partial_t\left[\frac{\mu y_t^2}{2}+\frac{T y_x^2}{2}\right]
+\partial_x(-T y_ty_x)=-\beta y_t^2.
\]
因此机械能密度 $u=\mu y_t^2/2+Ty_x^2/2$、能流 $\mathcal S=-Ty_ty_x$。按周期平均，利用复振幅公式
\[
\overline{\mathcal S}(x)=\frac{T\omega k_r A^2}{2}e^{-2k_ix}>0,
\quad
\overline{\mathcal P}_{\rm diss}/dx
=\frac{\beta\omega^2A^2}{2}e^{-2k_ix}.
\]
对 $x\in[0,\infty)$ 积分，$\overline P_{\rm diss}=\beta\omega^2A^2/(4k_i)$。由 $2k_rk_i=\beta\omega/T$ 得
\[
\overline{\mathcal S}(0)=\frac{T\omega k_rA^2}{2}
=\frac{\beta\omega^2A^2}{4k_i}
=\overline P_{\rm diss},
\]
输入功率完全转成黏滞耗散功率，没有假设波的平均能量流恒定。

(4) 若 $\beta/(\mu\omega)\ll1$，展开复平方根得
\[
k_r\simeq\omega\sqrt{\mu/T},\qquad
k_i\simeq\frac{\beta}{2\sqrt{\mu T}}.
\]
弱阻尼下相速约为 $\sqrt{T/\mu}$，而衰减长度对频率在首阶近似中不敏感。相反当 $\omega\ll\beta/\mu$ 时 $k^2\simeq i\beta\omega/T$，所以
\[
k_r\simeq k_i\simeq\sqrt{\frac{\beta\omega}{2T}}.
\]
这表现为低频黏滞扩散型响应，振幅穿透长度与相位尺度具有相同数量级。
